% !TeX program = xelatex
\documentclass{custhesis}

\custserial{}
\custdoctype{本科生毕业设计}
\custheadertitle{长春理工大学本科生毕业设计}
\custtitle{改进粒子群优化算法研究}
\custenglishtitle{Research on an Improved Particle Swarm Optimization Algorithm}
\custstudent{}
\custmajor{}
\custstudentid{}
\custsupervisor{}
\custschool{}
\custdate{二〇二六年六月}

\begin{document}

\makecustcover
\makecustdeclaration

\custfrontmatter

\begin{custcnabstract}
本文以智能优化算法研究为示例背景，围绕问题建模、基准算法分析、改进策略设计、实验方案和结果评价等环节，给出一套适用于计算机类算法研究方向本科毕业设计写作的示例性论文结构。示例内容以粒子群优化算法为基础，说明动态参数调整、搜索策略改进和实验对比分析的常见写法，并通过图表、公式和三线表展示常用论文元素的排版方式。

本文内容仅作为模板示例，重点展示长春理工大学计算机科学技术学院本科毕业设计中常见的版式、章节层级、图表编号、公式编号、参考文献和附录组织方式。实际使用时，应根据具体算法、数据集、实验结果和分析结论替换相应内容。
\custkeywords{粒子群优化\quad 启发式算法\quad 参数寻优\quad 收敛分析\quad 实验对比}
\end{custcnabstract}

\begin{custenabstract}
This template demonstrates a common thesis structure for undergraduate algorithm research in computer science. The sample topic is based on particle swarm optimization and covers problem modeling, baseline analysis, improvement strategy design, experimental setup, and result evaluation. It also provides examples of figures, tables, equations, references, appendices, and acknowledgements.

The sample text is intended only for formatting and structural reference. Users should replace the abstract, algorithm description, experimental results, and references with materials from their own projects before submission.
\custkeywordsen{particle swarm optimization; heuristic algorithm; parameter tuning; convergence analysis; experimental comparison}
\end{custenabstract}

\custtableofcontents

\custmainmatter

\chapter{绪论}

\section{研究背景}

优化问题广泛存在于路径规划、资源调度、参数寻优、模型训练和工程设计等场景中。许多实际问题具有非线性、多峰值、约束复杂或解析梯度难以获取等特点，传统确定性方法在求解效率、稳定性和适用范围上可能受到限制。启发式优化算法能够通过群体搜索、随机扰动和迭代更新等方式在复杂搜索空间中寻找较优解，是计算机类算法研究方向毕业设计中常见的选题类型。

算法研究类论文不仅需要说明算法流程是否能够运行，还应给出明确的问题定义、基准算法、改进动机、参数设置、实验数据和评价指标。论文写作中应重点阐明算法改进的依据、核心步骤、复杂度或收敛性讨论，以及与对比算法之间的实验差异。

\section{研究内容}

本文示例以改进粒子群优化算法为背景，展示算法研究类本科毕业设计正文中常见的内容组织方式。论文正文通常需要覆盖以下内容：

\begin{custenum}
  \item 建立待求解优化问题的数学模型，明确目标函数、约束条件和评价指标；
  \item 分析基准粒子群优化算法的基本流程、参数含义和可能存在的不足；
  \item 设计动态权重、局部扰动或种群更新等改进策略，并说明关键步骤；
  \item 设计对比实验，对收敛速度、求解精度、稳定性和运行时间进行分析。
\end{custenum}

\subsection{论文结构}

全文结构如下：第1章介绍课题背景、研究内容和论文组织结构；第2章说明问题建模、算法框架、改进策略和实验设计，并给出图表和公式排版示例；第3章总结全文工作并说明后续可改进方向。参考文献、附录和致谢按学校格式要求排版。

\chapter{算法设计与实验}

\section{问题建模}

以连续函数优化问题为例，设搜索空间为 $\Omega\subseteq\mathbb{R}^{d}$，算法目标是在给定迭代次数内寻找使目标函数取得较小值的候选解。算法研究类论文应明确问题输入、输出、约束条件、终止条件和评价指标。图\ref{fig:framework}给出了一个典型算法研究流程示例。

\begin{figure}[htbp]
  \centering
  \fbox{\parbox[c][4.2cm][c]{0.82\textwidth}{\centering 算法研究流程图占位\\问题建模、算法设计、参数设置、实验比较}}
  \caption{算法研究流程示意图}
  \label{fig:framework}
\end{figure}

\section{改进策略}

粒子群优化算法通过模拟群体协作搜索过程更新粒子位置。基础算法通常包含种群初始化、适应度计算、个体最优更新、全局最优更新和速度位置更新等步骤。为了提升算法在复杂多峰函数上的搜索能力，可引入动态惯性权重和局部扰动策略，在保持全局探索能力的同时增强后期局部开发能力。相关综述通常会从参数控制、拓扑结构和混合策略等角度分析算法改进方向\custcite{van2007particle}。

若论文中包含关键迭代公式、目标函数或评价指标，可使用公式环境。公式编号按章编排，如\custeqref{eq:update}所示：
\begin{equation}
  v_i^{t+1}=\omega_t v_i^t+c_1 r_1(p_i^t-x_i^t)+c_2 r_2(g^t-x_i^t)
  \label{eq:update}
\end{equation}
其中，$v_i^t$ 表示第 $i$ 个粒子在第 $t$ 次迭代时的速度，$x_i^t$ 表示粒子位置，$p_i^t$ 表示个体历史最优位置，$g^t$ 表示群体历史最优位置，$\omega_t$、$c_1$、$c_2$ 分别表示惯性权重和学习因子，$r_1$、$r_2$ 为随机数。

\section{实验设计与结果分析}

实验设计是算法研究类论文的重要组成部分。论文中应说明测试函数、数据集来源、参数设置、运行次数、评价指标和对比算法。表\ref{tab:datasets}展示了基准测试设置示例。

\begin{table}[H]
  \centering
  \caption{基准测试设置示例}
  \label{tab:datasets}
  \zihao{5}
  \begin{tabular}{cccc}
    \toprule
    测试函数 & 维数 & 搜索空间 & 评价指标 \\
    \midrule
    Sphere & 30 & $[-100,100]^d$ & 最优值、均值、标准差 \\
    Rastrigin & 30 & $[-5.12,5.12]^d$ & 最优值、均值、标准差 \\
    Rosenbrock & 30 & $[-30,30]^d$ & 最优值、均值、标准差 \\
    \bottomrule
  \end{tabular}
\end{table}

实验结果部分可从收敛曲线、最终适应度、稳定性和运行时间等角度展开分析。若结果表格较多，可将完整数据放入附录，正文保留关键对比结果和结论性讨论。根据参考规范，表题置于表上方，图题置于图下方，图表编号按章编排，表格一般采用三线表。

\chapter{结论}

本文给出了计算机专业本科毕业设计论文的一套通用写作结构和排版示例，覆盖封面、原创承诺书、中英文摘要、目录、正文、图表、公式、参考文献、附录和致谢等组成部分。正文示例围绕改进粒子群优化算法研究展开，展示了问题建模、算法流程、改进策略、实验设置和结果分析等算法类论文常见内容的组织方式。

实际写作时，应结合具体课题替换示例内容，补充真实的算法设计、参数设置、实验数据、对比结果和分析结论。后续可根据学院最新格式规范继续调整封面字段、章节样式、参考文献格式和附录组织方式。


\begin{custbibliography}
  \item 作者. 智能优化算法原理与应用[M]. 出版地: 出版社, 出版年: 起止页码.
  \item 作者. 改进粒子群优化算法研究[J]. 期刊名, 出版年份, 卷号(期号): 起止页码.
  \item Author A, Author B, Author C, et al. An improved heuristic optimization algorithm[J]. Journal Name, Year, Volume(Issue): pages.
  \item GB/T 7714-2015, 信息与文献 参考文献著录规则[S]. 北京: 中国标准出版社, 2015.
\end{custbibliography}

\custappendixstart
\custappendixchapter{实验设置}

\section{参数配置}

本附录用于记录算法实验环境、关键参数、补充图表或较长代码片段。附录中的图、表、公式编号与正文分开，使用 A、B、C 等字母作为前缀，例如图A.1、表A.1和式（A-1）。

\begin{table}[htbp]
  \centering
  \caption{算法实验环境与参数}
  \label{tab:appendix-env}
  \zihao{5}
  \begin{tabular}{cc}
    \toprule
    项目 & 配置 \\
    \midrule
    操作系统 & macOS / Linux \\
    编程语言 & Python 3.11 / MATLAB \\
    种群规模 & 30 \\
    最大迭代次数 & 500 \\
    独立运行次数 & 30 \\
    \bottomrule
  \end{tabular}
\end{table}


\begin{custacknowledgements}
这里填写致谢内容。致谢应简洁、真实，可对指导教师以及在毕业设计过程中给予帮助的组织和个人表示感谢。正式提交前请删除本段说明文字。
\end{custacknowledgements}

\end{document}
